\documentclass[11pt,a4paper]{article}

% Basic packages
\usepackage[utf8]{inputenc}
\usepackage[english]{babel}
\usepackage[margin=2.5cm]{geometry}
\usepackage{titlesec}
\usepackage{amsmath}
\usepackage{helvet}
\usepackage{setspace}

% Formatting adjustments to ensure it fits on one page
\titleformat{\section}{\normalfont\large\bfseries}{\thesection}{1em}{}[\vspace{-0.5em}\rule{\textwidth}{0.4pt}]
\titlespacing*{\section}{0pt}{3ex plus 1ex minus .2ex}{1.5ex plus .2ex}
\setlength{\parindent}{0pt}
\setlength{\parskip}{0.8em}

\begin{document}

\begin{center}
    \LARGE\textbf{Experimental Validation Proposal} \\
    \vspace{0.2cm}
    \large\textit{Assessment of a BBB-Permeable Cyclic Peptide Inhibitor of GSK3$\beta$}
\end{center}

\vspace{0.4cm}

\textbf{Context of the Computational Design:} This document outlines the experimental pipeline required to validate the \textit{in silico} design of a 15-amino acid peptide featuring a $\beta$-hairpin scaffold. This peptide has been computationally optimized to act as a competitive inhibitor of the ATP-binding pocket of the GSK3$\beta$ kinase. It is structurally stabilized by a specific intramolecular disulfide bond (Cys3--Cys13) and designed with favorable physicochemical properties to cross the blood-brain barrier (BBB), presenting a potential therapeutic strategy for Alzheimer's disease by mitigating Tau hyperphosphorylation.

\section*{Phase 1: Chemical Synthesis and Purification}
Given its short sequence length, the production of the peptide will be approached through chemical synthesis rather than recombinant expression. The peptide will be assembled using standard Fmoc Solid-Phase Peptide Synthesis (SPPS). Following chain assembly, a critical directed cyclization step will be performed under high-dilution conditions to strongly favor the formation of the intended intramolecular disulfide bond (Cys3--Cys13) while preventing intermolecular polymerization. Finally, the synthesized peptide will undergo rigorous quality control; it will be purified to $>95\%$ using preparative Reverse-Phase High-Performance Liquid Chromatography (RP-HPLC), and its exact mass and correct disulfide linkage will be confirmed via LC-MS or MALDI-TOF mass spectrometry.

\section*{Phase 2: Biophysical Characterization of Binding}
To empirically confirm the direct physical interaction between the synthetic peptide and its target, Surface Plasmon Resonance (SPR) will be employed. The recombinant kinase domain of GSK3$\beta$ will be immobilized on a sensor chip (e.g., CM5). Subsequently, serial injections of the peptide at varying concentrations will be flowed over the surface. This label-free, real-time analysis will allow for the precise determination of the equilibrium dissociation constant ($K_D$), as well as the kinetic association ($k_{on}$) and dissociation ($k_{off}$) rates, providing a comprehensive thermodynamic profile of the binding event.

\section*{Phase 3: \textit{In Vitro} Functional Assays}
Following the confirmation of binding, the inhibitory capacity and the specific mechanism of action will be evaluated through \textit{in vitro} kinase assays. Recombinant GSK3$\beta$, ATP, a specific peptide substrate, and the designed inhibitor will be incubated together. Luminescent assays (such as Kinase-Glo) will be utilized to quantify the residual ATP, serving as a proxy for kinase activity. Dose-response curves will be generated to calculate the half-maximal inhibitory concentration ($IC_{50}$). Furthermore, by conducting the assay under varying concentrations of ATP, Michaelis-Menten kinetics will be analyzed to definitively confirm that the peptide acts as a strictly competitive inhibitor at the ATP-binding site.

\section*{Phase 4: Blood-Brain Barrier Permeability Validation}
A crucial aspect of this project is the experimental validation of the predictions generated by the machine learning BBB classifiers developed computationally. Permeability will be assessed using two complementary models. Initially, a PAMPA-BBB (Parallel Artificial Membrane Permeability Assay) optimized with brain-specific lipids will be used to evaluate the peptide's capacity for passive transcellular diffusion. Subsequently, a more complex cellular model utilizing human brain microvascular endothelial cells (hCMEC/D3) cultured on polarized Transwell inserts will be employed. This cellular model allows for the assessment of both passive diffusion and active transport mechanisms (transcytosis), with the basolateral concentration of the peptide quantified via LC-MS/MS.

\section*{Phase 5: Cellular Efficacy and Cytotoxicity}
The final phase involves evaluating the therapeutic effect and safety profile of the peptide in a biologically relevant context. Neuronal cell lines (e.g., SH-SY5Y) will be treated with the peptide to assess its downstream impact on Tau protein hyperphosphorylation. Western Blot analyses utilizing phospho-specific antibodies (such as AT8) will be performed to quantify the ratio of phosphorylated Tau to total Tau protein. Concurrently, metabolic viability assays (e.g., MTT or resazurin) will be conducted to ensure that the peptide effectively inhibits the kinase without exhibiting intrinsic cytotoxicity at its therapeutically active concentrations.

\end{document}
